\writeSectionFrame{\presentationTitle}{Das Problem}
\frameWithImageOnly{\BILDERPFAD/problem1.jpg}
\note{
	Wir möchten gerne Komponenten für grafische Benutzeroberflächen entwickeln. Wir haben
	also Textfelder, Comboboxen, Listen, Schaltflächen usw. Sie müssen für den klassischen
	Desktop dargestellt werden.
}

\frameWithImageOnly{\BILDERPFAD/problem2.jpg}
\note{
	Wir möchten sie aber auch so dargestellt werden, dass bspw. Kinder damit arbeiten können.
	Die Schriften also größer, mehr Farben usw. 
}

\frameWithImageOnly{\BILDERPFAD/problem3.jpg}
\note{
	Wir benötigen eine weitere Variante für Smartphones.  
}

\begin{frame}{Mögliches Klassendiagramm}
    \begin{tikzpicture}[scale=0.7, transform shape]
        % Klassen
        \umlclass[x=0, y=0]{Component}{
            +paint() : void
        }{}
        
        \umlclass[x=-4, y=-3]{TextFieldPC}{}{
        }
        
        \umlclass[x=0, y=-3]{ComboBoxPC}{}{
        }
        
        \umlclass[x=4, y=-3]{ButtonPC}{}{
        }
        
        \umlclass[x=-4, y=2]{TextFieldSmartphone}{}{
        }
        
        \umlclass[x=4, y=2]{TextFieldChildren}{}{
        }

        \umlinherit{TextFieldPC}{Component}
        \umlinherit{ComboBoxPC}{Component}
        \umlinherit{ButtonPC}{Component}
        \umlinherit{TextFieldSmartphone}{Component}
        \umlinherit{TextFieldChildren}{Component}
    \end{tikzpicture}
\end{frame}

\note{
 	Wir könnten eine Oberklasse bauen, die für alle Elemente gleich ist. Sie könnte z.B. eine Methode
 	paint() enthalten, die alle Unterklassen entsprechend der ihnen zugrunde liegenden Bedingungen
 	oder Technologien implementieren müssen.
}

\begin{frame}{Wir haben ein Problem}
	\begin{center}
 		\myImageWithSource{0.45}{\BILDERPFAD/problem4.png}{https://refactoring.guru/design-patterns/bridge}
 	\end{center}
\end{frame}

\note{
 	Änderungen an einer solchen monolithischen Code-Basis sind immer sehr schwierig, weil man den
 	gesamten Monolithen verstehen muss. 
}